\ifdefined\gkver\else\def\gkver{student}\fi
\documentclass[\gkver]{gaokaozhenti}
\begin{document}

\chapter{2006年重庆理综}
%% paperType: 理综物理部分

\begin{enumerate}
\item
%% number: 14
%% paperName: 2006年重庆理综第 14 题 6 分
%% typeId: 1
%% type: 单选题
%% chapter: 曲线运动
%% point: 平抛运动
%% method: 无
%% score: 6
%% degree: 650
%% duplicateId: 0
%% body:
如右图，在同一竖直面内，小球 a、b 从高度不同的两点，分别以初速度 $v_{a}$ 和 $v_{b}$ 沿水平方向抛出，经过时间 $t_{a}$ 和 $t_{b}$ 后落到与两抛出点水平距离相等的 P 点。若不计空气阻力，下列关系式正确的是\xzanswer{A}

\onepicture[w=4.0cm]{\includesvg{figs/14}}

\fourchoices[answer=A]
{$t_{a}>t_{b}$，$v_{a}<v_{b}$}
{$t_{a}>t_{b}$，$v_{a}>v_{b}$}
{$t_{a}<t_{b}$，$v_{a}<v_{b}$}
{$t_{a}<t_{b}$，$v_{a}>v_{b}$}

%% answer: A
%% memo:
\memoanswer{由图可知，b 球落地的竖直高度小于 a 球的竖直高度，又 $\frac{1}{2}gt^{2}=h$，所以 $t_{a}>t_{b}$；两球落地的水平距离相等，由 $s=vt$ 可知，$s$ 相等，$t_{a}>t_{b}$，则 $v_{a}<v_{b}$，故 A 正确。}

\item
%% number: 15
%% paperName: 2006年重庆理综第 15 题 6 分
%% typeId: 1
%% type: 单选题
%% chapter: 曲线运动
%% point: 人造地球卫星
%% method: 无
%% score: 6
%% degree: 650
%% duplicateId: 0
%% body:
宇航员在月球上做自由落体实验，将某物体由距月球表面高 $h$ 处释放，经时间 $t$ 后落到月球表面（设月球半径为 $R$）。据上述信息推断，飞船在月球表面附近绕月球做匀速圆周运动所必须具有的速率为\xzanswer{B}

\fourchoices[answer=B]
{$\frac{2\sqrt{Rh}}{t}$}
{$\frac{\sqrt{2Rh}}{t}$}
{$\frac{\sqrt{Rh}}{t}$}
{$\frac{\sqrt{Rh}}{2t}$}

%% answer: B
%% memo:
\memoanswer{由 $G\frac{Mm}{R^{2}}=m\frac{v^{2}}{R}$ 得 $v=\sqrt{\frac{GM}{R}}$，由 $G\frac{Mm}{R^{2}}=mg$ 得 $GM=R^{2}g$，所以飞船绕月球表面做匀速圆周运动的速率为 $v=\sqrt{Rg}$；又在月球表面做自由落体实验时，从 $h$ 高释放，经 $t$ 秒落地，由 $\frac{1}{2}gt^{2}=h$ 得 $g=\frac{2h}{t^{2}}$，所以 $v=\frac{\sqrt{2Rh}}{t}$，故 B 正确。}

\item
%% number: 16
%% paperName: 2006年重庆理综第 16 题 6 分
%% typeId: 1
%% type: 单选题
%% chapter: 气体性质
%% point: 热力学第一定律
%% method: 无
%% score: 6
%% degree: 650
%% duplicateId: 0
%% body:
如右题图，某同学将空的薄金属筒开口向下压入水中。设水温均匀且恒定，筒内空气无泄漏，不计气体分子间相互作用，则被淹没的金属筒在缓慢下降过程中，筒内空气体积减小\xzanswer{C}

\onepicture[w=3.5cm]{figs/16.png}

\fourchoices[answer=C]
{从外界吸热}
{内能增大}
{向外界放热}
{内能减小}

%% answer: C
%% memo:
\memoanswer{金属筒在下降过程中，水对筒的压强增大，因此气体的压强增大，气体体积减小，外界对气体做功；又水温恒定，气体内能不变，而外界对气体做了功，则气体应向外界放热，故 C 正确。}

\item
%% number: 17
%% paperName: 2006年重庆理综第 17 题 6 分
%% typeId: 1
%% type: 单选题
%% chapter: 原子和原子核
%% point: 半衰期
%% method: 无
%% score: 6
%% degree: 650
%% duplicateId: 0
%% body:
$\ce{^{14}C}$ 是一种半衰期为 5730 年的放射性同位素，若考古工作者探测到某古木中 $\ce{^{14}C}$ 的含量为原来的 $\frac{1}{4}$，则该古树死亡时间距今大约\xzanswer{B}

\fourchoices[answer=B]
{22920 年}
{11460 年}
{5730 年}
{2865 年}

%% answer: B
%% memo:
\memoanswer{若某古树中 $\ce{^{14}C}$ 的含量为原来的 $\frac{1}{4}$，则该古树经历了 2 个半衰期，又 $\ce{^{14}C}$ 的半衰期为 5730 年，所以该古树死亡时间距今约为 11460 年，故 B 正确。}

\item
%% number: 18
%% paperName: 2006年重庆理综第 18 题 6 分
%% typeId: 1
%% type: 单选题
%% chapter: 机械波
%% point: 波的图象
%% method: 无
%% score: 6
%% degree: 650
%% duplicateId: 0
%% body:
右图为一列沿 $x$ 轴正方向传播的简谐横波在 $t=0$ 时的波形。当 R 点在 $t=0$ 时的振动状态传到 S 点时，PR 范围内（含 P、R）有一些质点正在向 $y$ 轴负方向运动，这些质点的 $x$ 坐标取值范围是\xzanswer{C}

\onepicture[w=7.0cm]{figs/18.png}

\fourchoices[answer=C]
{$2\Ucm\leq x\leq 4\Ucm$}
{$2\Ucm<x<4\Ucm$}
{$2\Ucm\leq x<3\Ucm$}
{$2\Ucm<x\leq 3\Ucm$}

%% answer: C
%% memo:
\memoanswer{当 R 点在 $t=0$ 时的振动状态传到 S 点时，波沿 $x$ 轴正方向传播的距离为 $\frac{3}{4}$ 个波长。将波形图向右平移 $\frac{3}{4}$ 个波长后，可以得到在 P、R 范围内，正在向 $y$ 轴负方向运动的质点的 $x$ 坐标取值范围是 $2\Ucm\leq x<3\Ucm$，故 C 正确。}

\item
%% number: 19
%% paperName: 2006年重庆理综第 19 题 6 分
%% typeId: 2
%% type: 多选题
%% chapter: 电场
%% point: 带电粒子的轨迹类问题
%% method: 无
%% score: 6
%% degree: 650
%% duplicateId: 0
%% body:
如右图，带正电的点电荷固定于 Q 点，电子在库仑力作用下，做以 Q 为焦点的椭圆运动。M、P、N 为椭圆上的三点，P 点是轨道上离 Q 最近的点。电子在从 M 经 P 到达 N 点的过程中\xzanswer{AC}

\onepicture[w=4.0cm]{\includesvg{figs/19}}

\fourchoices[answer=AC]
{速率先增大后减小}
{速率先减小后增大}
{电势能先减小后增大}
{电势能先增大后减小}

%% answer: AC
%% memo:
\memoanswer{Q 带正电，电子带负电，电子从 M 经 P 到达 N 点的过程中，电场力先做正功，后做负功，所以电子的动能先增大后减小，速率先增大后减小；又电场力做功等于电势能增量的负值，所以电势能先减小后增大，故 A、C 正确。}

\item
%% number: 20
%% paperName: 2006年重庆理综第 20 题 6 分
%% typeId: 2
%% type: 多选题
%% chapter: 光学
%% point: 光的折射
%% method: 无
%% score: 6
%% degree: 650
%% duplicateId: 0
%% body:
$\triangle$OMN 为玻璃等腰三棱镜的横截面。a、b 两束可见单色光从空气垂直射入棱镜底面 MN，在棱镜侧面 OM、ON 上反射和折射的情况如图所示。由此可知\xzanswer{BD}

\onepicture[w=4.5cm]{\includesvg{figs/20}}

\fourchoices[answer=BD]
{棱镜内 a 光的传播速度比 b 光的小}
{棱镜内 a 光的传播速度比 b 光的大}
{a 光的频率比 b 光的高}
{a 光的波长比 b 光的长}

%% answer: BD
%% memo:
\memoanswer{由图可知，a、b 两单色光以相同的入射角分别射向 OM、ON 两界面，b 光发生全反射，而 a 光发生折射，所以 a 光的折射率小于 b 光的折射率，因此 a 光的波长比 b 光的波长长；又 $v=\frac{c}{n}$，则在棱镜内 a 光的传播速度比 b 光的大，故 B、D 正确。}

\item
%% number: 21
%% paperName: 2006年重庆理综第 21 题 6 分
%% typeId: 2
%% type: 多选题
%% chapter: 电磁感应
%% point: 电磁感应受力分析
%% method: 无
%% score: 6
%% degree: 450
%% duplicateId: 0
%% body:
两根相距为 $L$ 的足够长的金属直角导轨如右图所示放置，它们各有一边在同一水平面内，另一边垂直于水平面。质量均为 $m$ 的金属细杆 ab、cd 与导轨垂直接触形成闭合回路，杆与导轨之间的动摩擦因数均为 $\mu$，导轨电阻不计，回路总电阻为 $2R$。整个装置处于磁感应强度大小为 $B$，方向竖直向上的匀强磁场中。当 ab 杆在平行于水平导轨的拉力 $F$ 作用下以速度 $v_{1}$ 沿导轨匀速运动时，cd 杆也正好以速度 $v_{2}$ 向下匀速运动。重力加速度为 $g$。以下说法正确的是\xzanswer{AD}

\onepicture[w=4.0cm]{\includesvg{figs/21}}

\fourchoices[answer=AD]
{ab 杆所受拉力 $F$ 的大小为 $\mu mg+\frac{B^{2}L^{2}v_{1}}{2R}$}
{cd 杆所受摩擦力为零}
{回路中的电流强度为 $\frac{BL(v_{1}+v_{2})}{2R}$}
{$\mu$ 与 $v_{1}$ 大小的关系为 $\mu=\frac{2Rmg}{B^{2}L^{2}v_{1}}$}

%% answer: AD
%% memo:
\memoanswer{由法拉第电磁感应定律可得，ab 杆在切割磁感线时产生的感应电动势为 $BLv_{1}$，cd 杆运动方向与磁场平行，无感应电动势，所以闭合回路中的感应电流 $I=\frac{E}{2R}=\frac{BLv_{1}}{2R}$。对 ab 杆，由平衡可知 $F=f+BIL$，所以 $F=\mu mg+\frac{B^{2}L^{2}v_{1}}{2R}$；同理对 cd 杆，由平衡可知 $f'=mg$，又 $f'=\mu BIL$，所以 $\mu=\frac{2Rmg}{B^{2}L^{2}v_{1}}$，故 A、D 正确。}

\item
%% number: 22
%% paperName: 2006年重庆理综第 22 题 13 分
%% typeId: 4
%% type: 实验题
%% chapter: 直线运动
%% point: 打点计时器公式
%% method: 无
%% score: 13
%% degree: 450
%% duplicateId: 0
%% body:
\begin{enumerate}
	\item
	用已调零且选择旋钮指向欧姆挡“$\times 10$”位置的多用电表测某电阻阻值，根据右图所示的表盘，被测电阻阻值为\tkanswer[2.0cm]{$220\UO$}。若将该表选择旋钮置于 $1\UmA$ 挡测电流，表盘仍如右图所示，则被测电流为\tkanswer[1.5cm]{$0.40\UmA$}。

	\onepicture[w=8.0cm, align=right]{figs/22a.png}

	\item
	某同学用如图（a）所示装置测量重力加速度 $g$，所用交流电频率为 $50\UHz$。在所选纸带上取某点为 0 号计数点，然后每 3 个点取一个计数点。所有测量数据及其标记符号如图（b）所示。

	该同学用两种方法处理数据（$T$ 为相邻两计数点的时间间隔）：方法 A：由 $g_{1}=\frac{s_{2}-s_{1}}{T^{2}}$，$g_{2}=\frac{s_{3}-s_{2}}{T^{2}}$，$\cdots\cdots$，$g_{5}=\frac{s_{6}-s_{5}}{T^{2}}$，取平均值 $g=8.667\Umsq$；方法 B：由 $g_{1}=\frac{s_{4}-s_{1}}{3T^{2}}$，$g_{2}=\frac{s_{5}-s_{2}}{3T^{2}}$，$g_{3}=\frac{s_{6}-s_{3}}{3T^{2}}$，取平均值 $g=8.673\Umsq$。从数据处理方法看，在 $s_{1}$、$s_{2}$、$s_{3}$、$s_{4}$、$s_{5}$、$s_{6}$ 中，对实验结果起作用的，方法 A 中有\tkanswer[2.0cm]{$s_{1}$、$s_{6}$}；方法 B 中有\tkanswer[3.8cm]{$s_{1}$、$s_{2}$、$s_{3}$、$s_{4}$、$s_{5}$、$s_{6}$}。因此，选择方法\tkanswer[1.2cm]{B}（A 或 B）更合理，这样可以减少实验的\tkanswer[1.5cm]{偶然}（系统或偶然）误差。本实验误差的主要来源有\tkanswer[3.5cm]{阻力（空气阻力，振针的阻力，限位孔的阻力，复写纸的阻力等），交流电频率波动，长度测量，数据处理方法等}（试举出两条）。

	\twopicture[align=right, showlabel=false, hA=4.0cm, hB=1.5cm]
	{figs/22b.png}
	{figs/22c.png}
\end{enumerate}

%% answer:
\jdanswer{
\begin{enumerate}
	\item
	$220\UO$；$0.40\UmA$

	\item
	$s_{1}$、$s_{6}$（或 $37.5$、$193.5$）；$s_{1}$、$s_{2}$、$s_{3}$、$s_{4}$、$s_{5}$、$s_{6}$（或 $37.5$、$69.0$、$100.5$、$131.5$、$163.0$、$193.5$）；B；偶然；阻力（空气阻力，振针的阻力，限位孔的阻力，复写纸的阻力等），交流电频率波动，长度测量，数据处理方法等。
\end{enumerate}
}

%% memo:
\memoanswer{
本题原卷及材料中未提供详解，待补充。
}

\item
%% number: 23
%% paperName: 2006年重庆理综第 23 题 16 分
%% typeId: 6
%% type: 计算题
%% chapter: 电路
%% point: 闭合电路欧姆定律
%% method: 无
%% score: 16
%% degree: 650
%% duplicateId: 0
%% body:
三只灯泡 $L_{1}$、$L_{2}$ 和 $L_{3}$ 的额定电压分别为 $1.5\UV$、$1.5\UV$ 和 $2.5\UV$，它们的额定电流都为 $0.3\UA$。若将它们连接成右图 1、图 2 所示电路，且灯泡都正常发光。

\begin{enumerate}
	\item
	试求图 1 电路的总电流和电阻 $R_{2}$ 消耗的电功率；
	\item
	分别计算两电路电源提供的电功率，并说明哪个电路更节能。
\end{enumerate}

\twopicture[align=right, showlabel=false, h=4.0cm]
{figs/23a.png}
{figs/23b.png}

%% answer:
\jdanswer{
\begin{enumerate}
	\item
	$I_{\text{总}}=0.9\UA$；$P_{R2}=0.045\UW$

	\item
	图 1 电源提供的电功率 $P_{\text{总}}=2.7\UW$；图 2 电源提供的电功率 $P'_{\text{总}}=1.8\UW$；由于灯泡都正常发光，两电路有用功率相等，而 $P'_{\text{总}}<P_{\text{总}}$，所以图 2 电路比图 1 电路节能。
\end{enumerate}
}

%% memo:
\memoanswer{
（1）由题意，在图 1 电路中，电路的总电流为 $I_{\text{总}}=I_{L1}+I_{L2}+I_{L3}=0.9\UA$，路端电压为 $U_{\text{路端}}=E-I_{\text{总}}r=2.55\UV$，则 $R_{2}$ 两端的电压与电流分别为 $U_{R2}=U_{\text{路端}}-U_{L3}=0.05\UV$，$I_{R2}=I_{\text{总}}=0.9\UA$，故电阻 $R_{2}$ 的消耗功率为 $P_{R2}=I_{R2}U_{R2}=0.045\UW$。

（2）图 1 电源提供的电功率为 $P_{\text{总}}=I_{\text{总}}E=0.9\times3\UW=2.7\UW$，图 2 电源提供的电功率为 $P'_{\text{总}}=I'_{\text{总}}E'=0.3\times6\UW=1.8\UW$，由于灯泡都正常发光，两电路有用功率相等，而 $P'_{\text{总}}<P_{\text{总}}$，所以图 2 电路比图 1 电路节能。
}

\item
%% number: 24
%% paperName: 2006年重庆理综第 24 题 19 分
%% typeId: 6
%% type: 计算题
%% chapter: 磁场
%% point: 带电粒子在磁场中的运动
%% method: 无
%% score: 19
%% degree: 450
%% duplicateId: 0
%% body:
有人设想用右图所示的装置来选择密度相同、大小不同的球状纳米粒子。粒子在电离室中电离后带正电，电量与其表面积成正比。电离后，粒子缓慢通过小孔 $O_{1}$ 进入极板间电压为 $U$ 的水平加速电场区域 Ⅰ，再通过小孔 $O_{2}$ 射入相互正交的恒定匀强电场、磁场区域 Ⅱ，其中磁场的磁感应强度大小为 $B$，方向如图。收集室的小孔 $O_{3}$ 与 $O_{1}$、$O_{2}$ 在同一条水平线上。半径为 $r_{0}$ 的粒子，其质量为 $m_{0}$、电量为 $q_{0}$，刚好能沿 $O_{1}O_{3}$ 直线射入收集室。不计纳米粒子重力。（$V_{\text{球}}=\frac{4}{3}\pi r^{3}$，$S_{\text{球}}=4\pi r^{2}$）

\begin{enumerate}
	\item
	试求图中区域 Ⅱ 的电场强度；
	\item
	试求半径为 $r$ 的粒子通过 $O_{2}$ 时的速率；
	\item
	讨论半径 $r\neq r_{0}$ 的粒子刚进入区域 Ⅱ 时向哪个极板偏转。
\end{enumerate}

\onepicture[w=8.0cm, align=right]{figs/24.png}

%% answer:
\jdanswer{
\begin{enumerate}
	\item
	$E=B\sqrt{\frac{2q_{0}U}{m_{0}}}$，方向竖直向上

	\item
	$v=\sqrt{\frac{r_{0}}{r}}v_{0}$

	\item
	$r>r_{0}$ 时，$v<v_{0}$，$F_{\text{合}}>0$，粒子会向上极板偏转；$r<r_{0}$ 时，$v>v_{0}$，$F_{\text{合}}<0$，粒子会向下极板偏转。
\end{enumerate}
}

%% memo:
\memoanswer{
（1）设半径为 $r_{0}$ 的粒子加速后的速度为 $v_{0}$，由动能定理得 $\frac{1}{2}m_{0}v_{0}^{2}=q_{0}U$，解得 $v_{0}=\sqrt{\frac{2q_{0}U}{m_{0}}}$。设区域 Ⅱ 内电场强度为 $E$，粒子在该区域做直线运动，由动能定理得 $q_{0}v_{0}B=q_{0}E$，解得 $E=v_{0}B=B\sqrt{\frac{2q_{0}U}{m_{0}}}$，电场强度方向竖直向上。

（2）设半径为 $r$ 的粒子的质量为 $m$、带电量为 $q$、被加速后的速度为 $v$，由题意知 $m=(\frac{r}{r_{0}})^{3}m_{0}$，$q=(\frac{r}{r_{0}})^{2}q_{0}$，由动能定理得 $\frac{1}{2}mv^{2}=qU$，解得 $v=\sqrt{\frac{2q_{0}Ur_{0}}{m_{0}r}}=\sqrt{\frac{r_{0}}{r}}v_{0}$。

（3）半径为 $r$ 的粒子，在刚进入区域 Ⅱ 时，受到的合力为 $F_{\text{合}}=qE-qvB=qB(v_{0}-v)$，由 $v=\sqrt{\frac{r_{0}}{r}}v_{0}$ 可知，当 $r>r_{0}$ 时，$v<v_{0}$，$F_{\text{合}}>0$，粒子会向上极板偏转；当 $r<r_{0}$ 时，$v>v_{0}$，$F_{\text{合}}<0$，粒子会向下极板偏转。
}

\item
%% number: 25
%% paperName: 2006年重庆理综第 25 题 20 分
%% typeId: 6
%% type: 计算题
%% chapter: 运动和力
%% point: 动量守恒定律
%% method: 无
%% score: 20
%% degree: 450
%% duplicateId: 0
%% body:
如右图，半径为 $R$ 的光滑圆形轨道固定在竖直面内。小球 A、B 质量分别为 $m$、$\beta m$（$\beta$ 为待定系数）。A 球从左边与圆心等高处由静止开始沿轨道下滑，与静止于轨道最低点的 B 球相撞，碰撞后 A、B 球能达到的最大高度均为 $\frac{R}{4}$，碰撞中无机械能损失。重力加速度为 $g$。试求：

\begin{enumerate}
	\item
	待定系数 $\beta$；
	\item
	第一次碰撞刚结束时小球 A、B 各自的速度和 B 球对轨道的压力；
	\item
	小球 A、B 在轨道最低处第二次碰撞刚结束时各自的速度，并讨论小球 A、B 在轨道最低处第 $n$ 次碰撞刚结束时各自的速度。
\end{enumerate}

\onepicture[w=5.0cm, align=right]{\includesvg{figs/25}}

%% answer:
\jdanswer{
\begin{enumerate}
	\item
	$\beta=3$

	\item
	$v_{A}=-\sqrt{\frac{gR}{2}}$，方向向左；$v_{B}=\sqrt{\frac{gR}{2}}$，方向向右；$N_{\text{压}}=-4.5mg$，方向竖直向下

	\item
	$v_{A}=-\sqrt{2gR}$；$v_{B}=0$；当 $n$ 为奇数时，小球 A、B 在第 $n$ 次碰撞刚结束时的速度分别与其第一次碰撞刚结束时相同；当 $n$ 为偶数时，小球 A、B 在第 $n$ 次碰撞刚结束时的速度分别与其第二次碰撞刚结束时相同。
\end{enumerate}
}

%% memo:
\memoanswer{
（1）由题意知，碰撞的过程中无机械能损失，由能量守恒得 $mgR=\frac{mgR}{4}+\frac{\beta mgR}{4}$，解得 $\beta=3$。

（2）设 A、B 碰撞后的速度分别为 $v_{1}$、$v_{2}$，由机械能守恒得 $\frac{1}{2}mv_{1}^{2}=\frac{mgR}{4}$，$\frac{1}{2}\beta mv_{2}^{2}=\frac{\beta mgR}{4}$，设向右为正、向左为负，解得 $v_{A}=-\sqrt{\frac{gR}{2}}$，即方向向左，$v_{B}=\sqrt{\frac{gR}{2}}$，即方向向右。设轨道对 B 球的支持力为 $N$，B 球对轨道的压力为 $N'$，方向竖直向上为正、向下为负，则 $N-\beta mg=\beta m\frac{v_{2}^{2}}{R}$，解得 $N'=-N=-4.5mg$，方向竖直向下。

（3）设 A、B 小球第二次碰撞刚结束时的速度分别为 $V_{1}$、$V_{2}$，由动量和能量守恒得 $mv_{1}-\beta mv_{2}=mV_{1}+\beta mV_{2}$，$mgR=\frac{1}{2}mV_{1}^{2}+\frac{1}{2}\beta mV_{2}^{2}$，解得 $v_{A}=-\sqrt{2gR}$，$v_{B}=0$。（另一组解 $V_{1}=2\sqrt{2gR}$，$V_{2}=-\sqrt{2gR}$，不合题意，舍去。）由此可得：当 $n$ 为奇数时，小球 A、B 在第 $n$ 次碰撞刚结束时的速度分别与其第一次碰撞刚结束时相同；当 $n$ 为偶数时，小球 A、B 在第 $n$ 次碰撞刚结束时的速度分别与其第二次碰撞刚结束时相同。
}

\end{enumerate}

\end{document}
